\honest{Leaky Generalizations}{Eliezer Yudkowsky}{2007}

Apples are usually good to eat; most people have ten fingers. Above the level of particles and fields, ``practically every generalization you use in the real world will be leaky.'' You just have to deal with it: if the cookie shop closes at 10 p.m. except at 6 p.m. on Thanksgiving, and today is Thanksgiving, come before six.

What gets in the way is ``need for closure'', the wish to settle a question once and for all. ``Raising the value of the stakes can increase need for closure'', just when we most need to tolerate complexity. \nb{The term is the psychologist Arie Kruglanski's. In Kruglanski's theory a cost of staying undecided, such as a deadline, raises the need for closure, but a cost of being wrong can raise a need to avoid closure. In the experiments Webster and Kruglanski report, fear of being evaluated reduced the effects that time pressure increased. The post's claim fits only stakes of the first kind.}

Life would be complicated even if what we wanted were simple, because rules for action inherit the leaks of the world. ``Instrumental values often have no specification which is both compact and local.'' Picture a box of money that a machine, worked with a dozen keys, opens or burns. There is more than one opening sequence, and too wrong a sequence burns the money. A short rule for which keys to press must describe the machine: press whatever opens the box. A rule about keys alone must be a giant table of every sequence. No rule is both short and about the keys alone. Worse, a rule may work almost always, such as pressing most keys three times, and miss the one key that burns the money if pressed once. You think you have a perfect rule when you have failed to picture all the machine's paths or to value all the side effects. The machine stands for the world, the box and the fire for our values, the keys for our actions. \nb{The box is built to have no simple rule, so it shows what the thesis means. That real instrumental values are ``often'' like this rests on the analogy.}

Given how many ways we value outcomes and how tangled the paths to them are, it is a wonder there is any helpful ethical advice at all, the strangest and still helpful being ``The end does not justify the means.'' But complicated actions do not show complicated goals. Some people smile wisely and say that morality is complicated: circumcision is right in one culture and wrong in another, torture is not always bad. But killing anyone, even Hitler, is ``a huge dose of negative terminal utility''. You should still shoot Hitler, because of all the lives that would be saved. \nb{The circumcision clause can be read as a claim that cultures hold different terminal values, which is how Yudkowsky, with ``I do not think'', had described that dispute a week earlier, in ``Terminal Values and Instrumental Values''. Read so, the Hitler reply does not address it.}

``Many'' infer from the good consequences of Hitler's death that his death itself is good, and so that ``Death is always a bad thing'' has exceptions. This is the type error of my earlier post, and a kind of double counting: it sets up a loop between expected utility and utility, where the flow should run one way, from utility to expected utility. Or it is the wish for a one-sided policy debate, where the best policy has no drawbacks. \nb{No one is named. The moral particularists, who deny that causing a death always counts the same way, argue from a general claim about reasons, not from consequences.} In my moral philosophy, the local disvalue of Hitler's death stays fixed whatever the consequences. \nb{This is close to W. D. Ross's prima facie duties (1930): a feature always counts, and an action with it can still be right overall.}

Punishing the evil might be argued good in itself, but not from the lives it saves: that shooting a man with a leveled gun saves others appeals to the value of life, not of death. Utilities might be leaky too. ``But it would be a separate argument.''
